\documentclass[11pt,a4paper]{article}
\usepackage[utf8]{inputenc}
\usepackage[T1]{fontenc}
\usepackage{geometry}
\geometry{top=1in, bottom=1in, left=1in, right=1in}
\usepackage{hyperref}
\usepackage{booktabs}
\usepackage{amsmath}
\usepackage{listings}
\usepackage{xcolor}
\usepackage{graphicx}
\usepackage{array}
\usepackage{fancyhdr}
\usepackage{titlesec}
\usepackage{longtable}

\hypersetup{
    colorlinks=true,
    linkcolor=purple,
    filecolor=magenta,      
    urlcolor=cyan,
    pdftitle=Genesis Engineering Specification (GES),
    pdfpagemode=FullScreen,
}

\definecolor{codegray}{rgb}{0.96,0.96,0.96}
\definecolor{codeborder}{rgb}{0.85,0.85,0.85}
\definecolor{commentgreen}{rgb}{0.12,0.55,0.22}
\definecolor{keywordblue}{rgb}{0.1,0.2,0.8}

\lstset{
    backgroundcolor=\color{codegray},
    basicstyle=\ttfamily\small,
    breakatwhitespace=false,
    breaklines=true,
    captionpos=b,
    commentstyle=\color{commentgreen},
    frame=single,
    rulecolor=\color{codeborder},
    keepspaces=true,
    keywordstyle=\color{keywordblue},
    language=Python,
    showspaces=false,
    showstringspaces=false,
    showtabs=false,
    tabsize=4
}

\pagestyle{fancy}
\fancyhf{}
\rhead{\small Genesis Engineering Specification (GES)}
\lhead{\small Project Genesis}
\cfoot{\thepage}

\title{\textbf{Genesis Engineering Specification (GES)}}
\author{Project Genesis Research Repository}
\date{\small Generated: July 2026}

\begin{document}

\maketitle
\thispagestyle{empty}

\newpage
\section*{Genesis Engineering Specification (GES)}
\subsection*{Version 1.0  APIs, Schemas, Database \& Deployment}

\textbackslash{}textbf{Project:} Project Genesis  Agentic Artificial Life \& Evolutionary Biology Simulator
\textbackslash{}textbf{Document Role:} Technical contracts, data schemas, infrastructure
\textbackslash{}textbf{Companion Documents:} GPS (Product)  GDS (Design System)  GCR (Content \& Research)
\textbackslash{}textbf{Status:} Living document  update with every API change or schema revision

\noindent\rule{\textwidth}{0.4pt}

\textgreater{} "The engineering specification is the contract between the simulation engine and the portal. If they disagree, the spec is the arbiter."

\noindent\rule{\textwidth}{0.4pt}

\subsection*{Part I  Repository Structure}

\subsubsection*{1.1 Why "portal" and not "website"}

The directory is named \textbackslash{}texttt{portal/}, not \textbackslash{}texttt{website/}.

"Website" describes the implementation (HTML served over HTTP).
"Portal" describes the purpose (a gateway to Genesis's scientific output).

Five years from now, the portal may include interactive simulations, community tools, downloadable datasets, and public APIs. The name must scale with the ambition.

\subsubsection*{1.2 Recommended Repository Layout}

\begin{lstlisting}[language=Python]
project-genesis/
│
├── engine/                          # Renamed from world/ — the simulation core
│   ├── state.py
│   ├── biomes.py
│   ├── climate.py
│   ├── climate_epochs.py
│   ├── erosion.py
│   ├── evolution_journal.py
│   ├── explainer.py
│   ├── generator.py
│   ├── habitability.py
│   ├── hydrology.py
│   ├── noise.py
│   ├── passability.py
│   ├── predictor.py
│   ├── resources.py
│   ├── telemetry.py
│   ├── world_presets.py
│   ├── checkpoint_io.py
│   └── agents/
│       ├── agent.py
│       ├── behavior_classifier.py
│       ├── cognitive.py
│       ├── decision.py
│       ├── drives.py
│       ├── genetics.py
│       ├── perception.py
│       ├── reproduction.py
│       └── simulation.py
│
├── portal/                          # The public research portal
│   ├── README.md
│   ├── GWAS.md                      # Complete spec (this document set)
│   └── ... (see §4 for full structure)
│
├── shared/                          # NEW — canonical contracts
│   ├── experiment_schema.json       # JSON Schema for valid experiment exports
│   └── types.ts                     # TypeScript types (auto-gen from schema)
│
├── tools/                           # NEW — CLI utilities
│   ├── export_experiment.py         # Validates + packages experiment as .zip
│   ├── validate_schema.py           # Validates folder against experiment_schema.json
│   └── generate_types.py            # Generates types.ts from experiment_schema.json
│
├── configs/                         # Unchanged
├── experiments/                     # Unchanged — raw local runs
├── research/                        # Unchanged — papers, notebooks
│
├── main.py
├── run.py
├── run_test.py
├── run_resume.py
├── run_research.py
└── visualizer.html
\end{lstlisting}

\subsubsection*{1.3 The \textbackslash{}texttt{shared/experiment\textbackslash{}\_schema.json} Contract}

This is the most important new file in the repository. It defines exactly what a valid Genesis experiment export must contain. The simulation engine validates against it on export. The portal validates against it on upload. If they disagree, the schema is correct.

See 3 (Canonical Experiment Schema) for the full schema definition.

\noindent\rule{\textwidth}{0.4pt}

\subsection*{Part II  Technology Stack}

\subsubsection*{2.1 Core Framework}

\begin{longtable}{lll}
\toprule
Layer & Choice & Rationale \\
\midrule
\endhead
Framework & \textbackslash{}textbf{Next.js 14} (App Router) & SSG for experiment pages (SEO + speed), ISR for cache invalidation on publish, API Routes for backend \\
Language & \textbackslash{}textbf{TypeScript} (strict) & Required for complex experiment data schemas. No \textbackslash{}texttt{any} types allowed. \\
Package Manager & \textbackslash{}textbf{pnpm} & Faster than npm, supports workspaces if needed \\
Styling & \textbackslash{}textbf{Vanilla CSS} (CSS Modules) & Full control over design system tokens. No TailwindCSS  the design system is too specific for utility classes. \\
Charts & \textbackslash{}textbf{Recharts} & Composable React chart library, excellent TypeScript support, supports custom rendering \\
Canvas (Replay, World Map) & \textbackslash{}textbf{Raw Canvas2D} & Full per-pixel control. No abstraction layer over rendering. \\
Lineage Tree & \textbackslash{}textbf{D3.js} (dendrogram only) & Use only for the hierarchical tree. Do not use D3 for anything else. \\
Schema Validation & \textbackslash{}textbf{Zod} & Runtime type-safe validation of experiment files on upload \\
ZIP Parsing & \textbackslash{}textbf{JSZip} & Client and server-side ZIP handling \\
CSV Parsing & \textbackslash{}textbf{Papa Parse} & Fast, streaming CSV parsing for large files \\
Markdown Rendering & \textbackslash{}textbf{react-markdown} + \textbackslash{}textbf{remark-gfm} & Render research documents with GFM table support \\
\bottomrule
\end{longtable}

\subsubsection*{2.2 Backend \& Infrastructure}

\begin{longtable}{lll}
\toprule
Layer & Choice & Rationale \\
\midrule
\endhead
Database & \textbackslash{}textbf{Supabase (Postgres)} & Managed Postgres + Auth + Storage in one platform. No separate backend needed. \\
Authentication & \textbackslash{}textbf{Supabase Auth} & Email/password for admin. Row-level security for all tables. \\
File Storage & \textbackslash{}textbf{Supabase Storage} & S3-compatible, CDN-backed, policy-based access control \\
Deployment & \textbackslash{}textbf{Vercel} & Zero-config Next.js deployment, edge functions, analytics \\
CDN & \textbackslash{}textbf{Cloudflare} & DNS, DDoS protection, caching in front of Vercel \\
Analytics & \textbackslash{}textbf{Vercel Analytics} & Privacy-friendly, built-in to Vercel plan \\
OG Images & \textbackslash{}textbf{Vercel OG (@vercel/og)} & Auto-generate per-experiment open graph images \\
\bottomrule
\end{longtable}

\subsubsection*{2.3 Development Tools}

\begin{longtable}{ll}
\toprule
Tool & Purpose \\
\midrule
\endhead
ESLint & Linting (use existing \textbackslash{}texttt{.eslintignore}) \\
Prettier & Formatting (use existing \textbackslash{}texttt{.prettierignore}) \\
TypeScript strict & No implicit any, strict null checks \\
\textbackslash{}texttt{next/image} & Optimized image delivery for world thumbnails \\
\bottomrule
\end{longtable}

\noindent\rule{\textwidth}{0.4pt}

\subsection*{Part III  Canonical Experiment Schema}

\subsubsection*{3.1 Required Files in Export ZIP}

An experiment \textbackslash{}texttt{.zip} is valid if and only if it contains all REQUIRED files and all files pass field validation.

\begin{longtable}{lll}
\toprule
File & Required & Description \\
\midrule
\endhead
\textbackslash{}texttt{config.json} & \textbackslash{}textbf{REQUIRED} & Full simulation config at save time \\
\textbackslash{}texttt{summary.json} & \textbackslash{}textbf{REQUIRED} & Derived metrics, distributions, leaderboards \\
\textbackslash{}texttt{events.json} & \textbackslash{}textbf{REQUIRED} & Chronological event timeline array \\
\textbackslash{}texttt{population.csv} & \textbackslash{}textbf{REQUIRED} & Tick-sampled population counts \\
\textbackslash{}texttt{timeline.csv} & \textbackslash{}textbf{REQUIRED} & Per-tick detailed stats \\
\textbackslash{}texttt{births.csv} & \textbackslash{}textbf{REQUIRED} & One row per birth event \\
\textbackslash{}texttt{deaths.csv} & \textbackslash{}textbf{REQUIRED} & One row per death event \\
\textbackslash{}texttt{genes.csv} & \textbackslash{}textbf{REQUIRED} & Per-agent gene values \\
\textbackslash{}texttt{colonies.csv} & \textbackslash{}textbf{REQUIRED} & Colony metadata \\
\textbackslash{}texttt{agent\textbackslash{}\_census.csv} & \textbackslash{}textbf{REQUIRED} & Per-agent final census \\
\textbackslash{}texttt{replay.json} & OPTIONAL & Full tick-by-tick replay data \\
\textbackslash{}texttt{checkpoint\textbackslash{}\_*.json} & OPTIONAL & Mid-run state checkpoints \\
\textbackslash{}texttt{evolution\textbackslash{}\_journal.md} & OPTIONAL & Auto-generated evolution journal \\
\textbackslash{}texttt{world.png} & OPTIONAL & World map thumbnail (512px) \\
\textbackslash{}texttt{report.md} & OPTIONAL & Auto-generated text summary report \\
\textbackslash{}texttt{simulation.log} & OPTIONAL & Full simulation log (download only) \\
\textbackslash{}texttt{performance\textbackslash{}\_log.csv} & OPTIONAL & Per-tick perf data \\
\textbackslash{}texttt{cognition.csv} & OPTIONAL & Per-agent cognitive stats \\
\textbackslash{}texttt{disputes\textbackslash{}\_history.json} & OPTIONAL & Territorial dispute records \\
\textbackslash{}texttt{dehydration\textbackslash{}\_death\textbackslash{}\_traces.json} & OPTIONAL & Dehydration telemetry \\
\textbackslash{}texttt{evolution\textbackslash{}\_replay.json} & OPTIONAL & Genetic history for replay slider \\
\bottomrule
\end{longtable}

\subsubsection*{3.2 \textbackslash{}texttt{config.json}  Full Field Specification}

\begin{lstlisting}[language=typescript]
interface ExperimentConfig {
  seed:                    number;     // World generation seed (integer)
  ticks:                   number;     // Total ticks completed (integer)
  scarcity:                number;     // Resource multiplier [0.1, 10.0]
  max_population:          number;     // Carrying capacity (integer)
  mutation_rate:           number;     // Gene mutation std dev [0.0, 0.2]
  reproduction_enabled:    boolean;
  disputes_enabled:        boolean;
  disasters_enabled:       boolean;
  healing_speed_mult:      number;
  shelter_build_speed_mult:number;
  shelter_search_dist:     number;
  planner_enabled:         boolean;
  sleep_consolidation_enabled: boolean;
  checkpoint_interval:     number;
  world_preset:            'island_chains' | 'arid_continent' | 'green_continent' |
                           'boreal_highlands' | 'tropical_ring' | null;
  climate_epoch_mode:      'legacy' | 'stable' | 'slow_change' | 'rapid_change' | 'random';
  ecology_ablation: {
    dehydration_ramp:      boolean;
    memory_fidelity:       boolean;
    water_caching:         boolean;
    deposit_utility_fix:   boolean;
  };
}
\end{lstlisting}

\subsubsection*{3.3 \textbackslash{}texttt{summary.json}  Full Field Specification}

\begin{lstlisting}[language=typescript]
interface ExperimentSummary {
  timestamp:      string;       // ISO 8601 datetime
  experiment:     string;       // Human name (from EXPERIMENT_NAME config var)
  seed:           number;
  ticks:          number;       // Actual ticks run (may differ from config if stopped early)
  scarcity:       number;
  survivors:      string;       // "0/18" format: "<alive>/<total_born>"
  avg_radius:     number;       // Average exploration radius in grid cells
  avg_discoveries:number;
  tests_passed:   boolean;
  max_generation: number;       // Highest generation number seen

  derived_metrics: {
    avg_generation_interval:     number;   // Ticks per generation
    population_doubling_time:    number;   // -1 if population never doubled
    colony_lifespans: {          [colonyName: string]: number };  // Ticks
    avg_genetic_diversity:       number;   // Shannon entropy [0.0, 1.0]
    food_efficiency:             number;
    water_efficiency:            number;
    energy_efficiency:           number;
    avg_prediction_error:        number;
    concept_formation_rate:      number;
    procedure_creation_rate:     number;
    avg_children:                number;
    avg_grandchildren:           number;
    avg_descendants:             number;
    avg_distance_from_birthplace:number;   // Grid cells
    avg_shelter_occupancy:       number;   // [0.0, 1.0]
    avg_social_degree:           number;   // Mean relationships per agent
    avg_cooperation_score:       number;
    avg_conflict_score:          number;
  };

  distributions: {
    lifespan:              { [bucket: string]: number };
    children:              { [bucket: string]: number };
    shelter:               { [level: string]: number };
    generation:            { [gen: string]: number };
    prediction_confidence: { [range: string]: number };
    concepts:              { [range: string]: number };
  };

  leaderboards: {
    longest_surviving: ColonyEntry[];
    largest_size:      ColonyEntry[];
    highest_health:    ColonyEntry[];
    most_generations:  ColonyEntry[];
    cognitive_mastery: ColonyEntry[];
  };
}

interface ColonyEntry {
  founder_id:              number;
  colony_id:               number;
  size:                    number;
  longevity_ticks:         number;
  max_generation:          number;
  avg_health:              number;
  avg_prediction_accuracy: number;
}
\end{lstlisting}

\subsubsection*{3.4 \textbackslash{}texttt{events.json}  Event Type Definitions}

\begin{lstlisting}[language=typescript]
type EventType = 'Birth' | 'Death' | 'Milestone' | 'Disaster' | 'Dispute' |
                 'ClimateEpoch' | 'Extinction';

interface BaseEvent {
  tick:        number;
  year:        number;
  day:         number;
  type:        EventType;
  description: string;
  metadata:    EventMetadata;
}

type EventMetadata =
  | BirthMetadata
  | DeathMetadata
  | MilestoneMetadata
  | DisasterMetadata
  | DisputeMetadata
  | ClimateEpochMetadata
  | ExtinctionMetadata;

interface BirthMetadata {
  child_id:     number;
  parent_a_id:  number;
  parent_b_id:  number;
  generation:   number;
  location:     [number, number];  // [y, x]
  colony_id:    number;
  colony_name:  string;
}

interface DeathMetadata {
  agent_id:    number;
  cause:       'dehydration' | 'starvation' | 'injury' | 'old_age' | 'cold' | 'heat' | 'unknown';
  age_ticks:   number;
  location:    [number, number];
  colony_id:   number;
}

interface MilestoneMetadata {
  population?: number;
  [key: string]: unknown;   // Free-form additional fields
}

interface DisasterMetadata {
  disaster_type: 'drought' | 'cold_wave' | 'heatwave' | 'famine';
  severity:      number;    // [0.0, 1.0]
}

interface DisputeMetadata {
  agent_a_id: number;
  agent_b_id: number;
  outcome:    'retreat_a' | 'retreat_b' | 'mutual_retreat' | 'fight';
  location:   [number, number];
}

interface ClimateEpochMetadata {
  epoch_name:        string;
  temperature_delta: number;    // °C offset
  rainfall_mult:     number;    // Multiplier on rainfall
}

interface ExtinctionMetadata {
  colony_id:      number;
  colony_name:    string;
  tick_founded:   number;
  tick_extinct:   number;
  total_members:  number;
}
\end{lstlisting}

\subsubsection*{3.5 CSV Column Headers}

\textbackslash{}textbf{\textbackslash{}texttt{population.csv}}
\begin{lstlisting}[language=Python]
tick,total,alpha,beta,gamma,delta
\end{lstlisting}

\textbackslash{}textbf{\textbackslash{}texttt{timeline.csv}}
\begin{lstlisting}[language=Python]
tick,year,day,alive,births_cumulative,deaths_cumulative,avg_health,
avg_hunger,avg_thirst,avg_fear,avg_generation,genetic_diversity
\end{lstlisting}

\textbackslash{}textbf{\textbackslash{}texttt{births.csv}}
\begin{lstlisting}[language=Python]
tick,year,day,child_id,parent_a_id,parent_b_id,generation,
colony_id,colony_name,location_y,location_x
\end{lstlisting}

\textbackslash{}textbf{\textbackslash{}texttt{deaths.csv}}
\begin{lstlisting}[language=Python]
tick,year,day,agent_id,cause,age_ticks,colony_id,
location_y,location_x
\end{lstlisting}

\textbackslash{}textbf{\textbackslash{}texttt{genes.csv}}
\begin{lstlisting}[language=Python]
agent_id,colony_id,generation,g_metabolism,g_thermoregulation,
g_vision,g_mobility,g_memory_fidelity,g_planning,g_novelty_seeking,
g_social_proximity,g_aggression,g_resource_sharing,g_risk_sensitivity,
g_resilience,g_longevity,g_learning_rate
\end{lstlisting}

\textbackslash{}textbf{\textbackslash{}texttt{agent\textbackslash{}\_census.csv}}
\begin{lstlisting}[language=Python]
agent_id,colony_id,colony_name,generation,age_ticks,lifespan_ticks,
health_at_death,children_count,shelter_level,cause_of_death,
birth_location_y,birth_location_x,death_location_y,death_location_x,
exploration_radius,discoveries
\end{lstlisting}

\textbackslash{}textbf{\textbackslash{}texttt{colonies.csv}}
\begin{lstlisting}[language=Python]
colony_id,colony_name,color_r,color_g,color_b,founder_count,
spawn_location_y,spawn_location_x,total_members,max_generation,
tick_founded,tick_extinct
\end{lstlisting}

\noindent\rule{\textwidth}{0.4pt}

\subsection*{Part IV  Database Schema (Supabase Postgres)}

\subsubsection*{4.1 Table: \textbackslash{}texttt{experiments}}

\begin{lstlisting}[language=sql]
CREATE TABLE experiments (
  -- Identity
  id              TEXT PRIMARY KEY,
  -- Format: 'EXP-YYYYMMDD-XXXX' (canonical) or timestamp slug
  -- Example: 'EXP-20260713-0001'

  slug            TEXT UNIQUE NOT NULL,
  -- URL-safe human title: 'territorial-conflict-extreme-scarcity'

  -- Metadata written by researcher on publish
  title           TEXT NOT NULL,
  abstract        TEXT,                        -- Researcher-written markdown
  tags            TEXT[] DEFAULT '{}',

  -- Key config fields (denormalized for fast filtering)
  seed            INTEGER NOT NULL,
  ticks           INTEGER NOT NULL,
  scarcity        REAL NOT NULL,
  world_preset    TEXT,
  climate_epoch_mode TEXT,
  max_population  INTEGER,
  mutation_rate   REAL,

  -- Key summary fields (denormalized for fast card rendering)
  total_agents    INTEGER,                     -- From survivors denominator
  survivors_count INTEGER,                     -- From survivors numerator
  max_generation  INTEGER,
  avg_genetic_diversity REAL,

  -- Full JSON blobs (for detail pages — not used in list queries)
  config_json     JSONB NOT NULL,
  summary_json    JSONB NOT NULL,

  -- Storage
  storage_path    TEXT NOT NULL,               -- 'experiments/{id}/' prefix
  thumbnail_url   TEXT,                        -- Public URL to world.png
  has_replay      BOOLEAN DEFAULT FALSE,       -- Whether replay.json is available
  replay_size_mb  REAL,                        -- File size for display

  -- Publishing
  is_published    BOOLEAN DEFAULT FALSE,
  is_featured     BOOLEAN DEFAULT FALSE,
  published_at    TIMESTAMPTZ,
  created_at      TIMESTAMPTZ DEFAULT NOW(),
  updated_at      TIMESTAMPTZ DEFAULT NOW()
);

-- Indexes for common filter/sort patterns
CREATE INDEX idx_experiments_published ON experiments(is_published, published_at DESC);
CREATE INDEX idx_experiments_scarcity  ON experiments(scarcity) WHERE is_published;
CREATE INDEX idx_experiments_preset    ON experiments(world_preset) WHERE is_published;
CREATE INDEX idx_experiments_featured  ON experiments(is_featured) WHERE is_published;
CREATE INDEX idx_experiments_tags      ON experiments USING GIN(tags);
\end{lstlisting}

\subsubsection*{4.2 Table: \textbackslash{}texttt{experiment\textbackslash{}\_events}}

\begin{lstlisting}[language=sql]
CREATE TABLE experiment_events (
  id              BIGSERIAL PRIMARY KEY,
  experiment_id   TEXT NOT NULL REFERENCES experiments(id) ON DELETE CASCADE,
  tick            INTEGER NOT NULL,
  year            INTEGER,
  day             INTEGER,
  event_type      TEXT NOT NULL,
  description     TEXT,
  metadata        JSONB
);

CREATE INDEX idx_events_exp_id   ON experiment_events(experiment_id);
CREATE INDEX idx_events_type     ON experiment_events(experiment_id, event_type);
CREATE INDEX idx_events_tick     ON experiment_events(experiment_id, tick);
\end{lstlisting}

\subsubsection*{4.3 Table: \textbackslash{}texttt{experiment\textbackslash{}\_agents}}

\begin{lstlisting}[language=sql]
CREATE TABLE experiment_agents (
  id              BIGSERIAL PRIMARY KEY,
  experiment_id   TEXT NOT NULL REFERENCES experiments(id) ON DELETE CASCADE,
  agent_id        INTEGER NOT NULL,
  colony_id       INTEGER,
  colony_name     TEXT,
  generation      INTEGER,
  age_ticks       INTEGER,
  lifespan_ticks  INTEGER,
  health_at_death REAL,
  children_count  INTEGER DEFAULT 0,
  shelter_level   INTEGER DEFAULT 0,
  cause_of_death  TEXT,
  birth_location  JSONB,    -- {y: number, x: number}
  death_location  JSONB,    -- {y: number, x: number}
  exploration_radius REAL,
  genes           JSONB     -- {g_metabolism: 0.72, g_vision: 0.44, ...}
);

CREATE INDEX idx_agents_exp_id   ON experiment_agents(experiment_id);
CREATE INDEX idx_agents_colony   ON experiment_agents(experiment_id, colony_id);
CREATE INDEX idx_agents_gen      ON experiment_agents(experiment_id, generation);
\end{lstlisting}

\subsubsection*{4.4 Table: \textbackslash{}texttt{experiment\textbackslash{}\_population}}

\begin{lstlisting}[language=sql]
CREATE TABLE experiment_population (
  id              BIGSERIAL PRIMARY KEY,
  experiment_id   TEXT NOT NULL REFERENCES experiments(id) ON DELETE CASCADE,
  tick            INTEGER NOT NULL,
  total           INTEGER,
  alpha           INTEGER,
  beta            INTEGER,
  gamma           INTEGER,
  delta           INTEGER
);

CREATE INDEX idx_pop_exp_tick ON experiment_population(experiment_id, tick);
\end{lstlisting}

\subsubsection*{4.5 Table: \textbackslash{}texttt{research\textbackslash{}\_documents}}

\begin{lstlisting}[language=sql]
CREATE TABLE research_documents (
  id              TEXT PRIMARY KEY,
  -- e.g. 'monograph', 'architecture', 'lab-notebook', 'thesis'

  title           TEXT NOT NULL,
  document_type   TEXT NOT NULL,
  -- 'monograph' | 'architecture' | 'lab_notebook' | 'thesis' | 'paper'

  content_md      TEXT NOT NULL,           -- Full markdown content
  storage_path    TEXT,                    -- PDF path in Supabase Storage
  version         TEXT,                    -- e.g. 'Phase 8.4'
  is_published    BOOLEAN DEFAULT FALSE,
  published_at    TIMESTAMPTZ,
  updated_at      TIMESTAMPTZ DEFAULT NOW(),
  created_at      TIMESTAMPTZ DEFAULT NOW()
);
\end{lstlisting}

\subsubsection*{4.6 Row-Level Security (RLS) Policies}

\begin{lstlisting}[language=sql]
-- Enable RLS on all tables
ALTER TABLE experiments       ENABLE ROW LEVEL SECURITY;
ALTER TABLE experiment_events ENABLE ROW LEVEL SECURITY;
ALTER TABLE experiment_agents ENABLE ROW LEVEL SECURITY;
ALTER TABLE experiment_population ENABLE ROW LEVEL SECURITY;
ALTER TABLE research_documents ENABLE ROW LEVEL SECURITY;

-- Public read access (published only)
CREATE POLICY "public_read_experiments"
  ON experiments FOR SELECT
  USING (is_published = TRUE);

CREATE POLICY "public_read_events"
  ON experiment_events FOR SELECT
  USING (EXISTS (
    SELECT 1 FROM experiments e
    WHERE e.id = experiment_id AND e.is_published = TRUE
  ));

-- Same pattern for agents, population, research_documents

-- Admin full access (authenticated user matching admin email)
CREATE POLICY "admin_all_experiments"
  ON experiments FOR ALL
  USING (auth.jwt() ->> 'email' = 'vinay@khosya.com');

-- Apply same admin policy to all tables
\end{lstlisting}

\subsubsection*{4.7 Supabase Storage Buckets}

\begin{longtable}{lll}
\toprule
Bucket & Path Convention & Access \\
\midrule
\endhead
\textbackslash{}texttt{experiment-zips} & \textbackslash{}texttt{{experiment\textbackslash{}\_id}/export.zip} & Private (admin write, signed URL read) \\
\textbackslash{}texttt{thumbnails} & \textbackslash{}texttt{{experiment\textbackslash{}\_id}/world.png} & Public \\
\textbackslash{}texttt{replays} & \textbackslash{}texttt{{experiment\textbackslash{}\_id}/replay.json} & Public (large files, CDN cached) \\
\textbackslash{}texttt{research-pdfs} & \textbackslash{}texttt{{doc\textbackslash{}\_id}/document.pdf} & Public \\
\textbackslash{}texttt{staging} & \textbackslash{}texttt{{uuid}/export.zip} & Private (temp, deleted after publish) \\
\bottomrule
\end{longtable}

\noindent\rule{\textwidth}{0.4pt}

\subsection*{Part V  REST API Specification}

All endpoints under \textbackslash{}texttt{/api/v1/}.

\subsubsection*{5.1 Experiments  Public}

\begin{lstlisting}[language=Python]
GET /api/v1/experiments
  Query: page=1, limit=12, worldPreset, minScarcity, maxScarcity,
         tags (comma-separated), sortBy=published_at|ticks|scarcity,
         sortOrder=desc|asc, featured=true
  Returns: {
    experiments: ExperimentCard[],
    total: number,
    page: number,
    totalPages: number
  }

GET /api/v1/experiments/:id
  Returns: ExperimentDetail

GET /api/v1/experiments/:id/events
  Query: types (comma-separated), page, limit
  Returns: { events: BaseEvent[], total: number }

GET /api/v1/experiments/:id/agents
  Query: colonyId, generation, sortBy=lifespan|generation|health
  Returns: { agents: AgentRecord[] }

GET /api/v1/experiments/:id/population
  Returns: { population: PopulationEntry[] }

GET /api/v1/experiments/:id/replay-url
  Returns: { url: string, expiresAt: string, sizeMb: number }
  Note: Returns a signed URL (valid 1 hour) to replay.json in Storage

GET /api/v1/experiments/compare
  Query: a (experiment_id), b (experiment_id)
  Returns: {
    a: ExperimentDetail,
    b: ExperimentDetail,
    diff: {
      [metricKey: string]: { a: number, b: number, delta: number, deltaPercent: number }
    }
  }

GET /api/v1/experiments/featured
  Returns: ExperimentCard  (the single featured experiment)
\end{lstlisting}

\subsubsection*{5.2 Research  Public}

\begin{lstlisting}[language=Python]
GET /api/v1/research
  Returns: { documents: ResearchDocumentSummary[] }

GET /api/v1/research/:id
  Returns: ResearchDocument (with full content_md)
\end{lstlisting}

\subsubsection*{5.3 Admin  Protected (requires Supabase Auth JWT)}

All admin routes validate the request has a valid JWT matching the admin account.

\begin{lstlisting}[language=Python]
POST /api/v1/admin/experiments/upload
  Body: multipart/form-data
    file: Blob (.zip)
  Returns: {
    stagingId: string,
    validation: ValidationReport,
    preview: ExperimentPreview  (auto-extracted metadata)
  }

POST /api/v1/admin/experiments/:stagingId/publish
  Body: {
    title:      string,
    abstract?:  string,    // Markdown
    tags:       string[],
    isFeatured: boolean,
    experimentId?: string  // Override auto-generated ID
  }
  Returns: { id: string, slug: string, url: string }

PATCH /api/v1/admin/experiments/:id
  Body: Partial<{title, abstract, tags, isFeatured, isPublished}>
  Returns: ExperimentDetail (updated)

DELETE /api/v1/admin/experiments/:id
  Returns: { success: boolean }

POST /api/v1/admin/research
  Body: {
    id:           string,
    title:        string,
    documentType: string,
    contentMd:    string,
    version?:     string
  }
  Returns: { id: string, url: string }

PATCH /api/v1/admin/research/:id
  Body: Partial<{title, contentMd, isPublished, version}>
  Returns: ResearchDocument
\end{lstlisting}

\subsubsection*{5.4 Response Type Definitions}

\begin{lstlisting}[language=typescript]
interface ExperimentCard {
  id:           string;
  slug:         string;
  title:        string;
  worldPreset:  string | null;
  scarcity:     number;
  ticks:        number;
  ticksHuman:   string;          // "7,970 ticks (22 yr)"
  totalAgents:  number;
  survivors:    number;
  survivorsStr: string;          // "0 of 18 survived"
  maxGeneration:number;
  publishedAt:  string;          // ISO 8601
  thumbnailUrl: string | null;
  isFeatured:   boolean;
  hasReplay:    boolean;
  tags:         string[];
}

interface ExperimentDetail extends ExperimentCard {
  abstract:     string | null;
  configJson:   ExperimentConfig;
  summaryJson:  ExperimentSummary;
  replayUrl?:   string;         // Only if hasReplay
}

interface ValidationReport {
  valid:    boolean;
  errors:   ValidationError[];
  warnings: ValidationWarning[];
  preview:  ExperimentPreview;
}

interface ValidationError {
  file:     string;             // 'summary.json'
  path:     string;             // '/derived_metrics/food_efficiency'
  expected: string;             // 'number'
  received: string;             // 'null'
  message:  string;             // Human-readable
}

interface ExperimentPreview {
  autoTitle:         string;    // From config experiment_name, formatted
  suggestedTags:     string[];  // ['island_chains', 'extreme_scarcity', ...]
  extractedMetrics: {
    seed, ticks, scarcity, worldPreset, totalAgents,
    survivorsStr, maxGeneration, colonyNames: string[]
  }
}
\end{lstlisting}

\noindent\rule{\textwidth}{0.4pt}

\subsection*{Part VI  Admin Upload \& Ingestion Pipeline}

\subsubsection*{6.1 Step-by-Step Pipeline}

\begin{lstlisting}[language=Python]
STEP 1 — RECEIVE
─────────────────────────────────────────────────────────────
User drops .zip onto upload zone.
Client: check file.type = 'application/zip' and size < 500MB.
Client: show upload progress bar.
Upload to: Supabase Storage bucket 'staging/{uuid}/export.zip'
On success: POST /api/v1/admin/experiments/upload with stagingId.

STEP 2 — EXTRACT
─────────────────────────────────────────────────────────────
Server: download zip from staging, extract in-memory (JSZip).
List all files in the archive.
Check all REQUIRED files exist.
If missing: return ValidationReport with specific missing files.

STEP 3 — PARSE & VALIDATE
─────────────────────────────────────────────────────────────
For each required file:

  config.json:
    Parse JSON. Validate against ExperimentConfig interface (Zod).
    Required: seed (integer), ticks (integer), scarcity (number > 0).

  summary.json:
    Parse JSON. Validate against ExperimentSummary interface (Zod).
    Cross-check: summary.seed must equal config.seed.
    Cross-check: summary.ticks must be <= config.ticks.

  events.json:
    Parse JSON. Must be an array.
    Validate each element against BaseEvent (Zod discriminated union on 'type').
    Warning (not error) if 0 events of a particular type (e.g. 0 disasters).

  population.csv (Papa Parse streaming):
    Required columns: tick, total, alpha, beta, gamma, delta.
    All values must be non-negative integers.
    Tick values must be monotonically increasing.

  timeline.csv, births.csv, deaths.csv, genes.csv, agent_census.csv:
    Required column headers checked.
    No data validation (too slow for large files — trusted from engine).

  Cross-validation:
    Every agent_id in births.csv must appear in genes.csv.
    Colony names in summary.json must match those in colonies.csv (if present).

STEP 4 — AUTO-EXTRACT METADATA
─────────────────────────────────────────────────────────────
From validated files, build ExperimentPreview:
  autoTitle:     Format config.experiment_name as title case
  suggestedTags: [world_preset, 'scarcity_extreme'|'scarcity_high'|'scarcity_normal',
                  climate_epoch_mode, 'fights_enabled' if disputes_enabled, ...]
  totalAgents:   Parse denominator from summary.survivors
  survivorsStr:  Format "X of Y survived"
  colonyNames:   From colonies.csv or default ['Alpha','Beta','Gamma','Delta']

STEP 5 — DATABASE INSERT (draft)
─────────────────────────────────────────────────────────────
Generate experiment_id: 'EXP-{YYYYMMDD}-{4-char-hash}'
Insert experiments row: is_published = FALSE
Batch-insert experiment_events rows (chunked at 1000 rows)
Batch-insert experiment_agents rows
Batch-insert experiment_population rows (downsampled to max 5000 rows)

STEP 6 — RETURN PREVIEW
─────────────────────────────────────────────────────────────
Return: { stagingId, validation, preview }
Client renders: full preview page (same layout as published page)
Client shows: metadata form (title, abstract, tags, featured toggle)

STEP 7 — PUBLISH
─────────────────────────────────────────────────────────────
POST /api/v1/admin/experiments/:stagingId/publish with form data.
Server:
  UPDATE experiments SET is_published=TRUE, published_at=NOW(), title, abstract, tags, isFeatured
  Copy zip from staging/ to experiments/{id}/export.zip
  Copy world.png (if present) to thumbnails/{id}/world.png
  Copy replay.json (if present) to replays/{id}/replay.json
  Update experiments.has_replay, thumbnail_url
  Invalidate Vercel cache: revalidatePath('/civilizations'), revalidatePath('/')
  Return: { id, slug, url }
Client: toast "Published — View live →"
\end{lstlisting}

\subsubsection*{6.2 Downsampling Strategy for \textbackslash{}texttt{experiment\textbackslash{}\_population}}

For experiments \textgreater{} 100,000 ticks, the raw population.csv may have \textgreater{}20,000 rows. Downsample to 5000 rows using:

\begin{lstlisting}[language=python]
# Adaptive downsampling: keep all rows for tick < 1000 (early dynamics),
# then sample at regular intervals for the remainder.
sample_every = max(1, len(rows) // 5000)
kept = [r for i, r in enumerate(rows) if r.tick < 1000 or i % sample_every == 0]
\end{lstlisting}

\subsubsection*{6.3 Replay File Strategy}

Replay files can be very large (a 1M-tick run would be enormous). Policy:

\begin{itemize}
  \item If \textbackslash{}texttt{replay.json} is \textless{} 50MB: upload as-is, served directly
  \item If \textbackslash{}texttt{replay.json} is 50200MB: upload compressed (gzip), decompress client-side
  \item If \textbackslash{}texttt{replay.json} is \textgreater{} 200MB: only upload first 50,000 ticks. Mark \textbackslash{}texttt{replay\textbackslash{}\_truncated: true} in experiments table.
  \item If \textbackslash{}texttt{replay.json} is absent: \textbackslash{}texttt{has\textbackslash{}\_replay = false}. Replay tab shows "Replay not available for this civilization."
\end{itemize}

\noindent\rule{\textwidth}{0.4pt}

\subsection*{Part VII  Folder Structure (Portal)}

\begin{lstlisting}[language=Python]
portal/
│
├── app/                                    # Next.js App Router
│   ├── page.tsx                            # / — Home
│   ├── layout.tsx                          # Root layout
│   │
│   ├── civilizations/
│   │   ├── page.tsx                        # /civilizations — Archive
│   │   ├── compare/
│   │   │   └── page.tsx                    # /civilizations/compare
│   │   └── [id]/
│   │       ├── page.tsx                    # /civilizations/[id] — Overview
│   │       ├── layout.tsx                  # Experiment layout (sidebar, tabs)
│   │       ├── observatory/page.tsx        # Charts & analytics
│   │       ├── replay/page.tsx             # Canvas replay
│   │       ├── chronicle/page.tsx          # Events timeline
│   │       ├── world/page.tsx              # World map
│   │       ├── lineage/page.tsx            # Agent roster + tree
│   │       ├── genome/page.tsx             # Gene distributions
│   │       └── conditions/page.tsx         # Config viewer
│   │
│   ├── research/
│   │   ├── page.tsx                        # Research hub
│   │   ├── monograph/page.tsx
│   │   ├── architecture/page.tsx
│   │   ├── field-notes/page.tsx
│   │   └── thesis/page.tsx
│   │
│   ├── timeline/page.tsx                   # Living Timeline
│   ├── origin/page.tsx                     # About Genesis
│   ├── collaborate/page.tsx                # Contact/Collaborate
│   ├── docs/
│   │   ├── page.tsx
│   │   ├── schema/page.tsx
│   │   ├── config/page.tsx
│   │   └── engine/page.tsx
│   │
│   ├── publications/page.tsx               # Reserved — empty state
│   ├── datasets/page.tsx                   # Reserved — empty state
│   ├── api/page.tsx                        # Reserved — empty state
│   ├── community/page.tsx                  # Reserved — empty state
│   ├── citations/page.tsx                  # How to cite Genesis
│   ├── version-history/page.tsx            # Engine changelog
│   │
│   └── admin/
│       ├── layout.tsx                      # Auth guard
│       ├── page.tsx                        # Command dashboard
│       ├── upload/page.tsx
│       ├── civilizations/page.tsx
│       └── research/page.tsx
│
├── api/                                    # Next.js API routes
│   └── v1/
│       ├── experiments/
│       │   ├── route.ts                    # GET /api/v1/experiments
│       │   ├── compare/route.ts
│       │   └── [id]/
│       │       ├── route.ts
│       │       ├── events/route.ts
│       │       ├── agents/route.ts
│       │       ├── population/route.ts
│       │       └── replay-url/route.ts
│       ├── research/
│       │   ├── route.ts
│       │   └── [id]/route.ts
│       └── admin/
│           ├── experiments/
│           │   ├── upload/route.ts
│           │   └── [stagingId]/publish/route.ts
│           └── research/route.ts
│
├── components/
│   ├── layout/
│   ├── civilization/
│   ├── charts/
│   ├── replay/
│   ├── world/
│   ├── agent/
│   ├── timeline/
│   ├── admin/
│   └── ui/
│
├── lib/
│   ├── supabase.ts
│   ├── supabase-server.ts                  # Server-side client (service role)
│   ├── api.ts
│   ├── schema-validator.ts                 # Zod schemas for all experiment files
│   ├── experiment-parser.ts               # ZIP extraction + CSV/JSON parsing
│   ├── experiment-ingestion.ts            # Full upload pipeline logic
│   └── replay-renderer.ts                 # Canvas2D rendering
│
├── hooks/
│   ├── useExperiment.ts
│   ├── useReplay.ts
│   └── useWorldMap.ts
│
├── types/
│   ├── experiment.ts
│   ├── agent.ts
│   ├── event.ts
│   └── genome.ts
│
├── styles/
│   ├── globals.css                         # Design system tokens (from GDS)
│   ├── typography.css
│   └── animations.css
│
├── public/
│   ├── fonts/
│   └── icons/
│
├── next.config.js
├── tsconfig.json
├── package.json
└── .env.local.example
\end{lstlisting}

\noindent\rule{\textwidth}{0.4pt}

\subsection*{Part VIII  Environment Variables}

\begin{lstlisting}[language=bash]
# .env.local.example

# Supabase (required)
NEXT_PUBLIC_SUPABASE_URL=https://your-project.supabase.co
NEXT_PUBLIC_SUPABASE_ANON_KEY=your-anon-key     # Safe to expose — RLS enforced
SUPABASE_SERVICE_ROLE_KEY=your-service-role-key  # NEVER expose to client

# Admin configuration
ADMIN_EMAIL=vinay@khosya.com                     # The only admin account

# Optional — Analytics
NEXT_PUBLIC_VERCEL_ANALYTICS_ID=your-id
\end{lstlisting}

\noindent\rule{\textwidth}{0.4pt}

\subsection*{Part IX  Deployment Architecture}

\subsubsection*{9.1 Infrastructure Diagram}

\begin{lstlisting}[language=Python]
User Browser
    │
    ▼ DNS lookup → Cloudflare (Proxy enabled)
Cloudflare
    │ CDN cache hit → return cached response
    │ cache miss →
    ▼
Vercel Edge Network
    ├─ Static assets (/fonts, /icons, /public) → CDN served
    ├─ SSG pages (/civilizations/[id]) → cached at build or ISR
    ├─ API routes (/api/v1/*) → Vercel serverless functions
    └─ Admin routes (/admin/*) → SSR, no caching
         │
         ▼
    Supabase (ap-south-1 or us-east-1)
         ├─ Postgres → experiments, events, agents, population, research
         ├─ Auth → admin session management
         └─ Storage → zips, thumbnails, replays, PDFs
\end{lstlisting}

\subsubsection*{9.2 Caching Strategy (ISR)}

\begin{lstlisting}[language=typescript]
// /app/civilizations/[id]/page.tsx
export const revalidate = 3600;    // Revalidate once per hour

// /app/civilizations/page.tsx
export const revalidate = 300;     // 5 minutes — new experiments show up quickly

// /app/research/[doc]/page.tsx
export const revalidate = 86400;   // 24 hours

// /app/page.tsx (homepage)
export const revalidate = 300;     // 5 minutes — featured experiment updates
\end{lstlisting}

On publish, the admin route calls:
\begin{lstlisting}[language=typescript]
import { revalidatePath } from 'next/cache';
revalidatePath('/civilizations');
revalidatePath('/');
revalidatePath(`/civilizations/${newId}`);
\end{lstlisting}

\subsubsection*{9.3 Build Configuration}

\begin{lstlisting}[language=javascript]
// next.config.js
const config = {
  images: {
    remotePatterns: [{
      protocol: 'https',
      hostname: '*.supabase.co',    // Allow Supabase Storage thumbnails
    }]
  },
  experimental: {
    serverComponentsExternalPackages: ['jszip', 'papaparse']
  }
};
\end{lstlisting}

\subsubsection*{9.4 Domain Setup}

\begin{lstlisting}[language=Python]
Primary: genesis.vinaykhosya.com
  → Cloudflare DNS (proxied)
  → Vercel project (custom domain configured)

Admin: genesis.vinaykhosya.com/admin
  → Same deployment, auth-guarded at middleware level
\end{lstlisting}

\noindent\rule{\textwidth}{0.4pt}

\subsection*{Part X  Security}

\subsubsection*{10.1 Admin Protection}

Every \textbackslash{}texttt{/admin/*} route is protected by Supabase Auth middleware:

\begin{lstlisting}[language=typescript]
// middleware.ts
import { createMiddlewareClient } from '@supabase/auth-helpers-nextjs'

export async function middleware(req: NextRequest) {
  const res = NextResponse.next();
  const supabase = createMiddlewareClient({ req, res });
  const { data: { session } } = await supabase.auth.getSession();

  if (req.nextUrl.pathname.startsWith('/admin')) {
    if (!session || session.user.email !== process.env.ADMIN_EMAIL) {
      return NextResponse.redirect(new URL('/login', req.url));
    }
  }
  return res;
}
\end{lstlisting}

\subsubsection*{10.2 Upload Security}

\begin{itemize}
  \item ZIP files are validated server-side before any database write
  \item Maximum file size: 500MB (enforced at Supabase Storage upload policy)
  \item Only \textbackslash{}texttt{.zip} MIME type accepted
  \item Staging files auto-deleted after 24 hours if not published (Supabase lifecycle rule)
\end{itemize}

\subsubsection*{10.3 API Rate Limiting}

\begin{itemize}
  \item Public endpoints: 100 requests/minute per IP (Vercel Edge Config)
  \item Admin endpoints: 20 requests/minute per session
  \item Upload endpoint: 5 requests/minute per session
\end{itemize}

\noindent\rule{\textwidth}{0.4pt}

\textbackslash{}textit{End of GES v1.0}

\noindent\rule{\textwidth}{0.4pt}
\textbackslash{}textbf{Document produced by Antigravity AI, July 2026}
\textbackslash{}textbf{Companion: GPS  GDS  GCR}

\end{document}
